\documentclass[10pt]{article}

% ---------- Packages ----------

\usepackage{parskip}
\usepackage{microtype}
\usepackage{xcolor}
\usepackage{geometry}
\geometry{margin=4em, top=8em,  headheight=\paperheight}
\usepackage[export]{adjustbox}
\usepackage{array}
\usepackage{amsmath}
\usepackage{amsfonts}
\usepackage{fancyhdr}
\usepackage{luacode}
\newcommand{\assessmentType}{Writing Response}
\newcommand{\topic}{}
\pagestyle{fancy}
\fancyhf{}
\lhead{Algebra II, \assessmentType\hspace{2em} First Name \rule{7em}{.4pt} Last Name \rule{7em}{.4pt}
}
\chead{}
\rhead{Period \rule{3em}{.4pt}}
\usepackage{lastpage}
\usepackage{xcolor}
\usepackage{enumitem}
\usepackage{pifont}
\usepackage{graphicx}
\graphicspath{{../img}}
\usepackage{pgfplots}
\pgfplotsset{compat=1.18}
\usepackage{tabularx}
\usepackage{polynom}
\usepackage{setspace}
\usepackage{pdfpages}
\usepackage{tikz}
\usetikzlibrary{angles,quotes}


% ---------- General formatting ----------
\setlength{\parindent}{0pt}
\setlength{\parskip}{3pt}
\setlist[enumerate]{
    leftmargin=1.8em,
    itemsep=3pt,
    topsep=2pt,
    parsep=1pt
}

% ---------- Ruled writing lines ----------
\newcommand{\writingline}{%
    \noindent
    \makebox[\linewidth][l]{%
        \rule{0pt}{1.65em}%
        \leaders\hrule height 0.35pt\hfill
    }\par
}

\newcommand{\R}{\mathbb R}
\newcommand{\e}{{\rm e}}
\newcommand{\pobr}[1]{\left\langle#1\right\rangle}
\newcommand{\norm}[1]{\lVert #1 \rVert}
\newcommand{\abs}[1]{\lvert #1 \rvert}

\DeclareMathOperator{\xd}{d\!}
\DeclareMathOperator{\proj}{proj}

\title{}
\date{}

\newcommand{\template}{
\noindent
{
\textbf{Tools:} scientific calculator \hspace{\stretch{1}} Score \rule{4em}{.4pt} 
}

\Huge
\begin{enumerate}[leftmargin=*, itemsep=1.2em]
\item
 The completely factored form of the expression \[x^6 - 63x^3 - 64\] is
 \begin{enumerate}[label = (\arabic*)]
\item
$(x-64)(x+1)$
\item
$(x^3 - 64)(x^3+1)$
\item
$(x+8)(x-8)(x+1)(x-1)$
\item
$(x-4)(x^2+4x+16)(x+1)(x^2-x+1)$
\end{enumerate} 
\end{enumerate}
\clearpage
}



\begin{document}

% ==========================================
% PAGE 1: Reading and reflection questions
% ==========================================

\begin{center}
    {\Large\bfseries Taking Ownership of Your Education}

    \vspace{2pt}
    {\small\itshape Reflecting on the First Marking Period}
\end{center}

\textbf{What does it mean to take ownership of your education?}

Education is more than attending classes, completing assignments,
and earning grades. It is an opportunity to develop knowledge,
skills, and ways of thinking that will serve you throughout your
life. Your teachers can guide you, explain difficult concepts,
and provide opportunities to learn. However, they cannot do the
learning for you. Ultimately, your education is something you
must actively participate in and take responsibility for.

\textbf{Taking ownership of your education means taking an active
role in your own learning.} It means recognizing that your choices,
habits, and actions matter. It also means being honest with yourself
about the challenges you face, recognizing what you can improve,
and taking steps toward your goals.

Taking ownership does not mean that everything is within your
control. You may face personal circumstances, family responsibilities,
health concerns, or other challenges that make attending class or
staying engaged difficult. You do not need to blame yourself for
circumstances beyond your control or disclose private information
to explain your experiences. Instead, focus on understanding your
situation, recognizing what you can influence, and seeking help
when you need it.

Ownership also means taking initiative when you encounter
difficulties. In mathematics, when you do not understand a concept,
you can ask questions, seek clarification, or explain the concept
to someone else to test your understanding. When you make a mistake,
you can examine your reasoning, identify what went wrong, and learn
from the experience instead of simply copying the correct answer.
These actions show that you are actively working to improve rather
than waiting for someone else to solve every difficulty for you.

Taking ownership also involves reflecting honestly on your habits
and decisions. If you miss classes, struggle to participate, or find
yourself disengaged, consider what contributed to these experiences.
Which factors were beyond your control? What choices could you make
differently? What resources or support might help you? The purpose
of reflection is not to judge yourself or make excuses, but to
understand your experiences and identify opportunities for growth.

The first marking period is an opportunity to reflect on how things
have gone so far. You cannot change the past, but you can learn from
it. As you look toward the next marking period, think about one or
two specific, realistic actions you can take to improve your
attendance, engagement, or learning habits. Rather than simply
promising to try harder, consider what you will actually do and
what support you may need.

\textbf{Your education is a shared effort, but your learning requires
your participation. You do not have to do everything alone, and
you do not have to be perfect. What matters is that you take an
active role in your own progress.}

\vspace{3pt}
\hrule
\vspace{4pt}

{\large\bfseries Written Reflection: My First Marking Period}

Reflect on your experiences during the first marking period and
write a thoughtful response addressing all three questions below.
Use specific examples where appropriate, but do not feel obligated
to disclose private information.

\begin{enumerate}
    \item \textbf{Attendance:} What were the major reasons for your
    absences during the first marking period? What circumstances
    or factors affected your ability to attend class regularly?

    \item \textbf{Classroom engagement:} What were the major reasons
    you sometimes found it difficult to participate in or stay
    engaged with class? Consider factors that affected your attention,
    participation, motivation, or willingness to ask for help.

    \item \textbf{Ownership and moving forward:} What does taking
    ownership of your education mean to you? Looking back on the
    first marking period, how have you taken ownership of your
    learning, and where could you improve? Identify one or two
    specific actions you will take during the next marking period
    to improve your attendance, classroom engagement, or learning
    habits. If you need support to make these changes, explain
    what kind of support would help.
\end{enumerate}

Be honest and thoughtful. The purpose of this assignment is not
to judge your past or assign blame, but to understand your
experiences, recognize your opportunities for growth, and make
a realistic plan for moving forward.
\clearpage

\noindent
{
\hspace{\stretch{1}} Score \rule{4em}{.4pt} 
}

\begin{center}
    {\Large\bfseries My Reflection}
\end{center}

% Use any remaining space on page 1 for writing.
\vfill
\writingline
\writingline
\writingline

\writingline
\writingline
\writingline
\writingline
\writingline
\writingline
\writingline
\writingline
\writingline
\writingline
\writingline
\writingline
\writingline
\writingline
\writingline
\writingline
\writingline
\writingline
\writingline
\writingline
\writingline
\writingline
\writingline
\writingline
\writingline
\writingline
\writingline
\writingline
\end{document}
